\documentclass[11pt]{article}
\usepackage[margin=1in]{geometry}
\usepackage{amsmath,amssymb,amsthm}
\newtheorem{theorem}{Theorem}
\newtheorem{lemma}[theorem]{Lemma}
\newtheorem{corollary}[theorem]{Corollary}
\newcommand{\N}{\mathcal N}
\newcommand{\C}{\mathcal C}
\title{Degree-Three Transversal Girth Hardness}
\author{NC0\_k-Avoid Laboratory V83}
\date{2026-08-02}
\begin{document}
\maketitle

\paragraph{Problem.}
A bipartite graph $B=(M,X;E)$ presents a transversal matroid on $M$.
The problem $\mathrm{D3\text{-}TRANSVERSAL\text{-}GIRTH}$ asks, under the
promise $|N_B(e)|\leq 3$ for every $e\in M$, whether the presented matroid has
a circuit of size at most a supplied threshold $L$.

\section{Path-selector series expansion}
For $e\in M$ order $N_B(e)=(x_{e,1},\ldots,x_{e,d_e})$.  Replace $e$ by
$C_e=\{e_1,\ldots,e_{d_e}\}$, introduce private right vertices
$P_e=\{p_{e,1},\ldots,p_{e,d_e-1}\}$, and put
\[
 N_{B'}(e_i)=\{x_{e,i}\}\cup
 \bigl(\{p_{e,i-1}\}:i>1\bigr)\cup
 \bigl(\{p_{e,i}\}:i<d_e\bigr).
\]
A source loop is retained as one loop.

\begin{lemma}[Private-chain matching]
Every proper subset of $C_e$ is matchable into $P_e$.
\end{lemma}
\begin{proof}
Choose a missing index $j$. Match $e_i$ to $p_{e,i}$ for $i<j$ and $e_i$ to
$p_{e,i-1}$ for $i>j$, then restrict this matching to the desired subset.
\end{proof}

\begin{lemma}[Complete-chain equivalence]
For every $A\subseteq M$, the set $A$ is matchable in $B$ if and only if
$C(A)=\bigcup_{e\in A}C_e$ is matchable in $B'$.
\end{lemma}
\begin{proof}
A source matching edge $e x_{e,j}$ is lifted by matching $e_j$ externally and
the remaining chain copies privately. Conversely, $C_e$ has $d_e$ left
vertices and only $d_e-1$ private right vertices, so every matching saturating
$C(A)$ uses at least one external neighbor for each chain. Selecting one such
edge per chain gives a matching of $A$ because the external vertices are
distinct.
\end{proof}

\begin{theorem}[Exact circuit correspondence]
The circuits of $B'$ are exactly
\[
 \left\{\bigcup_{e\in Q}C_e:Q\text{ is a circuit of }B\right\}.
\]
Hence the girth of $B'$ is the minimum source-circuit weight for
$w(e)=d_e$. If all $d_e=D$, then
\[
 \operatorname{girth}(T(B'))=D\operatorname{girth}(T(B)).
\]
Moreover, the maximum left degree of $B'$ is at most three.
\end{theorem}
\begin{proof}
A circuit of $B'$ cannot meet a chain in a nonempty proper subset: that part
has a private matching, so deleting it would leave a dependent proper subset,
contradicting minimality. Thus every circuit is a union of complete chains.
Complete-chain dependence and minimality are then equivalent to source
dependence and minimality by the preceding lemma. The girth statements follow.
\end{proof}

\section{Clique reduction}
The author-hosted final-version manuscript of the source theorem displays the
arithmetically inconsistent identity
$\ell={k\choose2}-k-1={k-1\choose2}$.  We do not use that identity.  Instead
we rederive the reservoir size required by Hall deficiency.

For a Clique instance $(G,k)$ with $k\geq4$, let
\[
 q={k\choose2},\qquad r=q-k-1,\qquad D=r+2=q-k+1.
\]
The source presentation has one left element for each edge $uv$ of $G$, right
side $V(G)\cup Z$ with $|Z|=r$, and neighborhood $\{u,v\}\cup Z$.

\begin{lemma}[Colbourn--Elmallah threshold]
Every edge set of size less than $q$ is matchable. An edge set of size $q$ is a
circuit if and only if it is the edge set of a $k$-clique.
\end{lemma}
\begin{proof}
For an edge set $H$ of size $t<q$ incident with $v$ graph vertices, Hall's
condition is $t-v\leq r$. If $v\geq k$, then
$t-v\leq(q-1)-k=r$. If $v\leq k-1$, simplicity gives
\[
 t-v\leq {v\choose2}-v\leq {k-1\choose2}-(k-1)\leq r.
\]
A $k$-clique has neighborhood size $r+k=q-1$ and is therefore a circuit.
Any other $q$-edge set uses at least $k+1$ graph vertices; its full
neighborhood has size at least $q$, while all proper subsets already satisfy
Hall's condition.
\end{proof}

\begin{theorem}[Degree-three transversal girth]
$\mathrm{D3\text{-}TRANSVERSAL\text{-}GIRTH}$ is NP-complete.
\end{theorem}
\begin{proof}
The problem belongs to NP because a dependent set of size at most $L$ is
verified by maximum matching. Apply the path-selector expansion to the
left-regular Clique presentation. Its maximum left degree becomes three and
its girth is multiplied by $D$. Consequently,
\[
 G\text{ has a }k\text{-clique}
 \quad\Longleftrightarrow\quad
 \operatorname{girth}(T(B'))\leq qD.
\]
The construction and threshold are polynomially bounded.
\end{proof}

\paragraph{Status boundary.}
The Clique template is from Colbourn and Elmallah, \emph{Discrete
Mathematics} 114 (1993), 115--129. The degree-reduction proof is internally
verified in V83. Novelty and peer-review status are not claimed. This theorem
does not resolve P versus NP or establish Boolean circuit lower bounds.

\end{document}
